\section{Carlo Michelstaedter}

\label{sec:Michelstaedter}

\begin{quotationx}
The is the part, in which \textbf{Julius Evola} details his debt to \textbf{Carlo Michelstaedter}. ``Conviction'', for Michelstaedter, is being and ``rhetoric'' is non-being, or becoming. Hence, he is attempting to describe the man who is convinced, i.e., the man who has being, who ``is''.

Opposed to him, is the man who is always lacking something. He looks to other people, conditions, or things, to fulfill himself. He imagines the future will somehow satisfy him, so he never lives in the present. You have known such types: they are always agitated, forever seeking novelty or excitement. They hide their sense of lack with rhetoric, creating entire worlds out of talk \emph{gerede} which they then inhabit. There is always an injustice, something to change, a new program to adopt, the politics of resentment, a glorious future that kills every current moment. This masks them from facing being in its totality. They don't even think with their heads, but with their tongues, as Carl Jung says.

Interestingly enough, the convinced man gives himself a mission of cosmic proportions, viz., to redeem the world.
\end{quotationx}


There is a man in whom the demand of the real individual toward absolute value, toward \emph{conviction}, has been confirmed in the modern epoch, like a lightning flash and in a reality intense with life; this man, who in the clearest way, by shattering all compromises by which the I, inadequate in itself, masks its \emph{abios bios} ``not living life'', Ariphon the Sicyonian , has been able to take life to its goal, compelling it to that which man has \emph{terror} of, more than any other thing in the world: to place oneself in the presence of oneself, to recognize oneself, to measure oneself at last with that scale which, alone, is the scale of \emph{value}, of \emph{being}---that man is \textbf{Carlo Michelstaedter}.

Previously, we alluded to some positions that however are forcefully asserted in Michelstaedter, and almost a tragedy, so as to see that his work to a great extent transcends the framework of an abstract discursive exposition. The fundamental point on which such a work is centered is the need for ``conviction'', i.e., of the I's absolute sufficiency in itself, understood as the \emph{real} principle of the individual. Now the concept of conviction is essentially characterized by Michelstaedter as the \emph{opposite of relationships}: that is, the I is not in itself, but in ``another'', it puts aside the principle of its own being, where its life is conditioned by things and relations, and there is not conviction, but rather lack, the death of value. Value is only what exists for itself, it does not ask from anyone the principle of its own life and its own power---autarchy. So not only the whole of life consisting of needs, feelings, social conventions, intellectual embellishments, etc., but even the corporeal organism itself and the system of nature (that is, understood as generated, in its infinitely recurrent spatial-temporal development, from the interminable gravitation with which lack pursues being that, however, insofar as it seeks it outside itself, will never succeed in possessing it), is brought back into the sphere of non-value.

The I that is convinced that it is as \emph{continuous} as it keeps itself outside the fullness of actual possession and pushes its conviction onto a later moment, by which it is made dependent; the I flees from itself in every moment, it does not have itself but seeks itself and \emph{desires} itself, and yet it will never be able to have that in any future, since the future is the very symbol of its privation, the shadow that follows the man who flees from the body of his reality that is maintained at every unchanged point. This is the meaning of everyday life for Michelstaedter except, in one thing, non-value, that which must not be. As opposed to such a situation, this is the voice of conviction: to endure, to resist lack with all one's life at every point, not to give up on life---which is lacking in itself, by looking outside or in the future---not to demand, but to hold
\emph{being} in one's own fist: \emph{not to become, but to endure}.


\begin{quotex}
While lack, ``always anxious about the future, hastens time, and replaces one empty present with the next, the stability of the individual occupies infinite time in the present and stops time. Its firmness is a vertiginous path for all others who are in the current. Its every moment is a century in the life of others---as long as he sets himself on fire and succeeds in enduring into the utmost present''.
\end{quotex}


To identify such a point, it is rather important to understand the nature of the connection that is contained in the premises: since the world is understood to be produced from the direction of lack, whose concrete incarnation it almost is, (and in that, the idealistic demand remains satisfied, in fact, it is reinforced in significance, because the genesis of the real, as in Buddhism, is connected to a moment of value, to a direction of will) it is an illusion to think that the point of conviction can be realized through an abstract inner and subjective endurance in a value that, as in the Stoic, would have as opposed to being (nature) such that it is, while not having value.

Those who ask for conviction must instead rise to a world responsibility, the work of conviction is essentially cosmic. I must not flee my lack --- or the world --- but take it on myself, adapt myself to its load and redeem it. Michelstaedter in fact says: ``You cannot be called convinced as long as something still is, which is not conviction'' and points to conviction as the ``extreme consciousness of someone who is one with things, has in himself all things: \emph{en sunexes} tr. continuous ''. The concrete point of conviction would therefore have the sense of a cosmic consummation.



In order to illuminate Michelstaedter's central problem, it may be useful to connect the concept of insufficiency or lack to the Aristotelian concept of the imperfect act. The imperfect or impure act is that of powers that do not achieve actuality from themselves, but require the involvement of the other. An example is the act of sensible perception; in it the power of perceiving is not sufficient in itself, it does not produce perception from itself, but needs correlation to the object. Now the fundamental point, which is linked to Michelstaedter's position, is this: the imperfect act resolves the privation of the I only apparently, since it in reality reconfirms it.

For example, the I is thirsty: as long as it drinks, it will continue to be thirsty, since by drinking he will prove the point of anyone who is not sufficient to his own life, but who needs the ``other'' in order to live, who \emph{is} not but \emph{becomes}: water and the rest are only the symbols of his deficiency (it is necessary to pay attention to this: \emph{one does not desire because there is a privation of being, but there is a privation of being because one desires} and, in the second place: \emph{there is only a desire}, e.g. thirst, \emph{because there are certain things}, e.g. water, \emph{but the things desired, just like the privation of the being that drives toward it, are created at birth by the desire for a relative, which therefore is the prius that creates the correlation as its two terms}, viz., privation and the related object, thirst and water), and since one feeds on it and \emph{demands} life from it, the I feeds itself only from its own privation and remains in it, fleeing from that pure or perfect act, from that eternal water, of which Christ speaks (John 4:14-15), so that every thirst, just as every other privation, will be forever defeated. This craving, this obscure drive that brings the I toward the outer world --- toward the other --- is that which generates the system of finite and contingent reality. Conviction, which burns at the point of absolute endurance, therefore, such a drive of the pure being-in-itself also has the sense of a consummation of the world.

Now it must be understood in the sense of such a consummation. Here various consequences arise that Michelstaedter did not completely treat. First of all, to say that I must not flee my lack means, among other things, to say that I must recognize myself as the creative function of the world: the justification of idealism follows from it (i.e., the system that says the world is posited by the I) according to a moral imperative. But the world, according to the premise, is recognized as the denial of value. From the general imperative of redeeming the world, of accepting the person of lack, then proceeds, always as the moral postulate, that is, not as theoretical construct, but rather as the object of a morally imperative affirmation --- of a practice. A second point is that the same denial of value \emph{must} be recognized in a certain way as a value. This point is important. If in fact I mean the desire that has generated the world as a brute given, as an irrational absolute, it is evident that conviction, conceived as its denial, \emph{depends} on it, therefore it is not in itself unconditionally sufficient but rather owes its life to another, that lets affirm itself in denying itself. In such a case, i.e., in the case that the same craving must not be taken back to an order included in the affirmation of value but remains absolutely a given, conviction would therefore not be at all conviction --- the initial mystery would inexorably damn its perfection into an illusion.

Therefore, as a moral postulate it is necessary to admit that the same antithesis participates in a certain way in value. But in what way? Such a question brings to life the concept of conviction into a dynamic principle. It is clear in fact that if conviction sharpens itself to a pure, unrelated sufficiency --- i.e., to a state --- rather than to sufficiency as denial of an insufficiency --- i.e. to an \emph{act}, to a \emph{relation} --- the antithesis certainly has a value and is explained: in a first moment, the I \emph{must} posit privation, a non-value, even if under the condition that it is posited solely because it must be denied, since this act of negation, and this alone, makes the value of conviction come to light.

Besides, what does it mean to deny the antithesis --- which here is the same as saying nature? If you recall that for Michelstaedter nature is a nonvalue as a symbol, the incarnation of the flight of the I from actual possession of itself, since it is correlative to an imperfect or impure act. It is not therefore a question of annulling this or that determination, since by doing that, one would only hit the effect, not the transcendental root of the non-value; not even annulling in general every action, since the antithesis is not action qua action in general, but rather action insofar as it flees from itself, insofar as ``becoming'', and it is not that every action is necessarily such. Rather, what is necessary to resolve is the passive, heteronomous, extraverted mode of action. Now the negation of such a mode is constituted by sufficient action, or actions on the basis of power, which was previously explained. To live every act on the basis of perfect possession is therefore to transfigure the whole of universal determination to the point of expressing only the same body of infinite \emph{potestas}, of the absolute Individual made of power, such is therefore the meaning of cosmic consummation. As the concreteness of ``rhetoric'' is the development of the world of dependency and necessity, so the concreteness of persuasion is the development of the world of autarchy as the cosmic dominator, and the point of the austere denial is only the neutral point (\emph{laya}, ``dissolution, extinction'') between the two phases.

The development of Michelstaedter's views in magical idealism thus results pursuant to logical continuity. On the other hand, Michelstaedter remained in a certain way closed to indeterminate negation and this, in great part, in order not to be fixed sufficiently on that, which, finite and infinite, is not an object or a particular action, but rather a particular way of living whatever object or action. The true Master does not, in general, need to deny (in the sense of annulling) and, with the pretense of rendering it absolute, to exalt life into an immobile undifferentiated unity or sudden inspiration: the creative act, the act of power --- that is not an act of desire or violence, but an act of a gift --- rather than destroying perfect possession, it witnesses to it and reconfirms it. It is that Michelstaedter, captivated in the immediacy of the terrible intensity in which he experienced the exigency of absolute value, unable to give this a concrete form and therefore to deploy it in the doctrine of power: we can perhaps attribute the tragic end of his mortal existence to that.


\flrightit{Posted on 2014-06-13 by Cologero}

\begin{center}* * *\end{center}

\begin{footnotesize}
\begin{sffamily}
\texttt{Logres on 2015-04-18 at 21:21 said:}

``Rather, what is necessary to resolve is the passive, heteronomous, extraverted mode of action.''

This sounds a lot like ``making every task play'' or making ``all yokes easy''.

``The true Master does not, in general, need to deny (in the sense of annulling) and, with the pretense of rendering it absolute, to exalt life into an immobile undifferentiated unity or sudden inspiration: the creative act, the act of power ---that is not an act of desire or violence, but an act of a gift---rather than destroying perfect possession, it witnesses to it and reconfirms it. ''

The first sentence here, ``to exalt life'' perfectly describes the Romantic movement in poetry (not to deny them their insights), at least historically, as well as my own early tendency in life. This is perhaps the finest sentence I've read so far in Evola.

\hfill

\texttt{Logres on 2015-04-18 at 21:26 said:}

So much of modern life seems to revolve around a confused and inchoate grasp of ``the cloud and the tree are One, to the eye of the child''. They are, but not as dissolved in a mush to be possessed and gobbled as part of the mass, but as the energies of the same All Creative God, which (I learn, thanks to Exile and Cologero) even have different qualities (creating, sustaining, glorifying). One, but never the same. Never the same, truly Infinite, always One. Logos.

\end{sffamily}
\end{footnotesize}